\chapter[Prime-index encodings and the weighted bound]{Prime-index encodings\\and the weighted bound}\label{ch:weighted}
\sourcebox{erdos1060\_complete\_partial\_bound.tex}{Complete working proof of the $C_0$ bound. The chain-packing bound in the next chapter is numerically stronger, but the injective encodings here remain reusable.}
\section{Statement and relation to the conjecture}
Throughout, $\sigma(k)=\sum_{d\mid k}d$, $h(k)=k\sigma(k)$, and $v_p(k)$ is the exponent of the prime $p$ in $k$. All logarithms are natural. We use empty products with value $1$.

The weaker conjecture in the formulation of Erd\H{o}s Problem 1060 recorded in \cite{kominers} is
\begin{equation}\label{weighted:eq:conjecture}
\log\max\{1,f(N)\}=o\!\left(\frac{\log N}{\log\log N}\right).
\end{equation}
In particular, it requires the following bound for \emph{every} $\varepsilon>0$, with a threshold allowed to depend on $\varepsilon$:
\[
f(N)\le\exp\!\left(\varepsilon\frac{\log N}{\log\log N}\right)
\quad\text{for all sufficiently large }N.
\]
The stronger proposed bound is $f(N)\le(\log N)^{O(1)}$.

The established result proved here is the following.
\begin{theorem}\label{weighted:thm:main}
Set
\[
C_0=\frac{\log3}{2}-\frac{\log2}{3}=0.3182570841474065\ldots.
\]
For every integer $N\ge2$ and every real $z\ge3$, writing $\alpha_p=v_p(N)$, we have
\begin{equation}\label{weighted:eq:finite}
\boxed{\log\max\{1,f(N)\}\le
\sum_{\substack{p\mid N\\p\le z}}\log(\alpha_p+1)
+C_0\frac{\log N}{\log z}.}
\end{equation}
Consequently, uniformly as $N\to\infty$,
\begin{equation}\label{weighted:eq:asymptotic}
\boxed{\log\max\{1,f(N)\}\le
(C_0+o(1))\frac{\log N}{\log\log N}.}
\end{equation}
\end{theorem}

Kominers proves $f(N)\le\prod_{p\mid N}v_p(N)$ and its consequence with constant $\log3/3$ in place of $C_0$ \cite[Theorem 1.2 and Corollary 1.3]{kominers}. We do not use that theorem as a premise. The argument below is a self-contained weighted-counting presentation of the prime-index bound in the preceding working note \cite{internal}.


\section{An injective encoding for each prime index}
For a prime $\ell$ and a positive integer $t$, define the $\ell$-free part
\[
t_{\ell'}=\frac{t}{\ell^{v_\ell(t)}}.
\]

\begin{lemma}\label{weighted:lem:encoding}
Fix a prime $\ell$. Suppose that $h(k)=h(m)$, that $v_\ell(k)=v_\ell(m)$, and that
\[
(v_p(k)+1)_{\ell'}=(v_p(m)+1)_{\ell'}
\qquad\text{for every prime }p\ne\ell.
\]
Then $k=m$.
\end{lemma}

\begin{proof}
The function $h$ is multiplicative: if $\gcd(u,v)=1$, then $h(uv)=h(u)h(v)$. In the equality between the prime-power product formulas for $h(k)$ and $h(m)$, cancel the equal local factors.

At any remaining prime $p$, denote the larger exponent by $a$, the smaller by $b$, and put $d=a-b>0$ and $t=b+1$. The equality of the labels gives
\[
a+1=t\ell^j\qquad\text{for some }j\ge1.
\]
Therefore
\begin{equation}\label{weighted:eq:ratio}
\frac{h(p^a)}{h(p^b)}=p^dQ_p,
\qquad Q_p=\frac{p^{t\ell^j}-1}{p^t-1}
=\sum_{i=0}^{\ell^j-1}p^{it}\in\NN.
\end{equation}

Every prime $q$ dividing $Q_p$ satisfies
\begin{equation}\label{weighted:eq:congruence}
q=\ell\quad\text{or}\quad q\equiv1\pmod\ell.
\end{equation}
To prove this, first note that $q\ne p$, since $Q_p\equiv1\pmod p$. If $p^t\equiv1\pmod q$, the sum in \eqref{weighted:eq:ratio} gives $0\equiv\ell^j\pmod q$, so $q=\ell$. Otherwise the order $r=\ord_q(p)$ divides $t\ell^j$ but does not divide $t$. All prime factors of $r$ other than $\ell$ occur to exponents at most their exponents in $t$. Thus $v_\ell(r)>v_\ell(t)$, so $\ell\mid r\mid q-1$. This proves \eqref{weighted:eq:congruence}.

Let $P_+$ be the primes at which $k$ has the larger exponent and $P_-$ the primes at which $m$ has the larger exponent. Define
\[
A=\prod_{p\in P_+}p^{d_p},\qquad
B=\prod_{p\in P_-}p^{d_p},\qquad
U=\prod_{p\in P_+}Q_p,\qquad
V=\prod_{p\in P_-}Q_p.
\]
Then $AU=BV$ and $\gcd(A,B)=1$. Hence $A\mid V$ and $B\mid U$. Every differing input prime divides an opposite-side factor $Q_q$. The exponent at $\ell$ was fixed, so \eqref{weighted:eq:congruence} implies
\begin{equation}\label{weighted:eq:inputcongruence}
p\equiv1\pmod\ell\qquad(p\in P_+\cup P_-).
\end{equation}

For each such $p$, the sum in \eqref{weighted:eq:ratio} shows that $\ell\mid Q_p$. Also
\[
Q_p=p^{d_p}\frac{1-p^{-(d_p+t)}}{1-p^{-t}}
<\frac{p^{d_p}}{1-p^{-t}}
\le\frac p{p-1}p^{d_p}.
\]
Since $p\ge\ell+1$,
\[
\frac{p}{\ell(p-1)}\le\frac{\ell+1}{\ell^2}<1.
\]
It follows that
\begin{equation}\label{weighted:eq:strict}
(Q_p)_{\ell'}\le\frac{Q_p}{\ell}<p^{d_p}.
\end{equation}

The integers $A,B$ are coprime to $\ell$. Taking $\ell$-free parts in $AU=BV$ and using $\gcd(A,B)=1$ gives
\[
A\mid\prod_{p\in P_-}(Q_p)_{\ell'},\qquad
B\mid\prod_{p\in P_+}(Q_p)_{\ell'}.
\]
If any exponent differs, multiplication and \eqref{weighted:eq:strict} yield
\[
AB\le\prod_{p\in P_+\cup P_-}(Q_p)_{\ell'}
<\prod_{p\in P_+\cup P_-}p^{d_p}=AB,
\]
a contradiction. Hence no exponent differs and $k=m$.
\end{proof}


\section{Local weights for the two encodings}
For an integer $a\ge0$, put
\[
\CC_\ell(a)=\{r\in\NN:1\le r\le a+1,\ \ell\nmid r\}.
\]
This is exactly the set of labels $(e+1)_{\ell'}$ as $e$ ranges over $0,\ldots,a$. Write
\[
c=\frac12\log\frac32,\qquad \lambda=\frac13\log2.
\]

\begin{lemma}\label{weighted:lem:weights}
For every $a\ge0$, there are positive weights $w_{\ell,a}(r)$, for $\ell=2,3$ and $r\in\CC_\ell(a)$, such that
\begin{equation}\label{weighted:eq:normalized}
\sum_{r\in\CC_\ell(a)}w_{\ell,a}(r)=1\qquad(\ell=2,3),
\end{equation}
and for every $e\in\{0,\ldots,a\}$,
\begin{equation}\label{weighted:eq:local}
\boxed{w_{2,a}((e+1)_{2'})^{2/3}
       w_{3,a}((e+1)_{3'})^{1/3}
       \ge e^{-ca-\lambda e}.}
\end{equation}
\end{lemma}

\begin{proof}
When $a=0$, assign weight $1$ to the sole label for each index.

When $a=1$, set
\[
w_{2,1}(1)=1,\qquad w_{3,1}(1)=\frac23,\qquad w_{3,1}(2)=\frac13.
\]
The left side of \eqref{weighted:eq:local} equals $(2/3)^{1/3}2^{-e/3}$, which is at least $(2/3)^{1/2}2^{-e/3}=e^{-c-\lambda e}$.

When $a=2$, set
\[
w_{2,2}(1)=\frac23,\quad w_{2,2}(3)=\frac13,
\qquad
w_{3,2}(1)=\frac23,\quad w_{3,2}(2)=\frac13.
\]
For $e=0,1,2$, respectively, the pair of labels is $(1,1)$, $(1,2)$, $(3,1)$. Direct substitution gives equality in \eqref{weighted:eq:local} in all three cases:
\[
w_{2,2}((e+1)_{2'})^{2/3}w_{3,2}((e+1)_{3'})^{1/3}
=\frac23\,2^{-e/3}=e^{-2c-\lambda e}.
\]

Now let $a\ge3$. Put $t=2^{-1/3}$ and
\[
Z_\ell(a)=\sum_{r\in\CC_\ell(a)}t^{r-1},\qquad
w_{\ell,a}(r)=\frac{t^{r-1}}{Z_\ell(a)}.
\]
The normalization \eqref{weighted:eq:normalized} is immediate. We will verify
\begin{equation}\label{weighted:eq:partition}
Z_2(a)^2Z_3(a)\le\left(\frac32\right)^{3a/2}.
\end{equation}
If $r_\ell=(e+1)_{\ell'}$, then $r_\ell-1\le e$. Since $0<t<1$, \eqref{weighted:eq:partition} gives
\begin{align*}
w_{2,a}(r_2)^{2/3}w_{3,a}(r_3)^{1/3}
&=\frac{t^{\frac23(r_2-1)+\frac13(r_3-1)}}{Z_2(a)^{2/3}Z_3(a)^{1/3}}\\
&\ge\frac{t^e}{(3/2)^{a/2}}
=e^{-ca-\lambda e},
\end{align*}
as required.

It remains to prove \eqref{weighted:eq:partition}. The elementary inequality $t<4/5$ follows from $(4/5)^3>1/2$. Also $t^3=1/2$.

For $a=3$,
\[
Z_2=1+t^2<\frac{41}{25},\qquad
Z_3=1+t+t^3<\frac{23}{10}.
\]
The exact rational comparison
\begin{equation}\label{weighted:eq:check3}
\left[\left(\frac{41}{25}\right)^2\frac{23}{10}\right]^2
<\left(\frac32\right)^9
\end{equation}
proves \eqref{weighted:eq:partition}.

For $a=4,5$,
\[
Z_2=1+t^2+t^4<\frac{1281}{625},\qquad
Z_3=\frac32+t+t^4<\frac{3387}{1250}.
\]
Here
\begin{equation}\label{weighted:eq:check45}
\left(\frac{1281}{625}\right)^2\frac{3387}{1250}
<\left(\frac32\right)^6,
\end{equation}
which proves \eqref{weighted:eq:partition} for $a=4$ and hence also for $a=5$.

Finally, for $a\ge6$, summing infinite geometric series gives
\[
Z_2(a)\le\frac1{1-t^2}<\frac{25}{9},\qquad
Z_3(a)\le\frac{1+t}{1-t^3}=2(1+t)<\frac{18}{5}.
\]
Thus
\begin{equation}\label{weighted:eq:check6}
Z_2(a)^2Z_3(a)<\left(\frac{25}{9}\right)^2\frac{18}{5}
=\frac{250}{9}<\left(\frac32\right)^9
\le\left(\frac32\right)^{3a/2}.
\end{equation}
For explicit verification of the rational checks, the right side minus the left side in \eqref{weighted:eq:check3} and \eqref{weighted:eq:check45}, respectively, is
\[
\frac{878868043}{5000000000}>0,
\qquad
\frac{124598601}{15625000000}>0.
\]
Also $(3/2)^9-250/9=49147/4608>0$. No numerical approximation is a premise of these inequalities.
\end{proof}


\section{Weighted counting and the global bound}
\begin{proof}[Proof of Theorem~\ref{weighted:thm:main}]
If $f(N)=0$, \eqref{weighted:eq:finite} is immediate. Otherwise partition all preimages of $N$ according to their exact exponents at the primes $p\le z$. Every preimage divides $N$, so the number of parts is at most
\[
T=\prod_{\substack{p\mid N\\p\le z}}(\alpha_p+1).
\]
Fix a nonempty part $\FF$. For $\ell=2,3$ and $k\in\FF$, define
\[
W_\ell(k)=\prod_{\substack{p\mid N\\p>z}}
 w_{\ell,\alpha_p}\bigl((v_p(k)+1)_{\ell'}\bigr).
\]
The exponents at $2$ and $3$ are fixed because $z\ge3$. By Lemma~\ref{weighted:lem:encoding}, each large-prime label vector, for either index $\ell=2$ or $\ell=3$, is injective on $\FF$. Thus its values are distinct points in the Cartesian product of the corresponding label sets. Using \eqref{weighted:eq:normalized},
\begin{equation}\label{weighted:eq:sumweights}
\sum_{k\in\FF}W_\ell(k)
\le\prod_{\substack{p\mid N\\p>z}}
   \sum_{r\in\CC_\ell(\alpha_p)}w_{\ell,\alpha_p}(r)
=1.
\end{equation}

By the weighted arithmetic--geometric mean inequality,
\[
W_2(k)^{2/3}W_3(k)^{1/3}
\le\frac23W_2(k)+\frac13W_3(k).
\]
Summing and applying \eqref{weighted:eq:sumweights} gives
\begin{equation}\label{weighted:eq:sumupper}
\sum_{k\in\FF}W_2(k)^{2/3}W_3(k)^{1/3}\le1.
\end{equation}
On the other hand, multiplying \eqref{weighted:eq:local} over the large primes yields
\[
W_2(k)^{2/3}W_3(k)^{1/3}
\ge\exp\!\left(-c\sum_{\substack{p\mid N\\p>z}}\alpha_p
                 -\lambda\sum_{\substack{p\mid N\\p>z}}v_p(k)\right).
\]
We have
\[
\sum_{\substack{p\mid N\\p>z}}\alpha_p
\le\frac{\log N}{\log z}.
\]
Also $\sigma(k)\ge k$, so $k^2\le h(k)=N$ and
\[
\sum_{\substack{p\mid N\\p>z}}v_p(k)
\le\frac{\log k}{\log z}
\le\frac{\log N}{2\log z}.
\]
Hence, uniformly for $k\in\FF$,
\[
W_2(k)^{2/3}W_3(k)^{1/3}
\ge\exp\!\left(-\left(c+\frac\lambda2\right)
                       \frac{\log N}{\log z}\right)
=\exp\!\left(-C_0\frac{\log N}{\log z}\right).
\]
Together with \eqref{weighted:eq:sumupper}, this gives
\[
|\FF|\le\exp\!\left(C_0\frac{\log N}{\log z}\right).
\]
There are at most $T$ parts. Multiplying by $T$ and taking logarithms proves \eqref{weighted:eq:finite}.

For \eqref{weighted:eq:asymptotic}, write $L=\log N$, $u=\log L$, and choose
\[
z=\frac{L}{u^3},
\]
which is at least $3$ for all sufficiently large $N$. Since $\alpha_p\le L/\log2$ and there are at most $z$ primes not exceeding $z$,
\[
\sum_{\substack{p\mid N\\p\le z}}\log(\alpha_p+1)
\le z\log\!\left(\frac L{\log2}+1\right)
=O\!\left(\frac L{u^2}\right)
=o\!\left(\frac L u\right).
\]
Moreover,
\[
\log z=u-3\log u\sim u.
\]
Substitution into \eqref{weighted:eq:finite} gives
\[
\log\max\{1,f(N)\}
\le(C_0+o(1))\frac L u
=(C_0+o(1))\frac{\log N}{\log\log N},
\]
completing the proof. No prime number theorem is needed.
\end{proof}


\section{Other consequences and the local entropy barrier}
For each prime index $\ell$, counting the injective labels gives
\[
 f(N)\le(\alpha_\ell+1)\prod_{p\mid N,\,p\ne\ell}
 \left(\alpha_p+1-\left\lfloor\frac{\alpha_p+1}{\ell}\right\rfloor\right).
\]
For $\ell=2$ this is
\[
 f(N)\le(\alpha_2+1)\prod_{p\mid N,\ p\text{ odd}}
 \left(\left\lfloor\frac{\alpha_p}{2}\right\rfloor+1\right).
\]
Using $\log(\lfloor a/2\rfloor+1)\le a\log2/2$ and the usual prime split
proves the older unrestricted constant $\log2/2$. That majorant itself has
constant $\log2/2$ along products of odd prime squares.

The equivalent entropy formulation of the local-weight proof is
\[
 \frac23H((J+1)_{2'})+\frac13H((J+1)_{3'})
 \le\frac a2\log(3/2)+\frac{\log2}{3}\,\mathbb EJ,
 \qquad 0\le J\le a.
\]
Here entropy uses natural logarithms. It follows from the local weights by
$H(Y)\le-\mathbb E\log w(Y)$ and averaging the two inequalities. For a uniform
$J\in\{0,1,2\}$, both label entropies equal
$\log3-\tfrac23\log2=2C_0$, while $\mathbb EJ=1=a/2$.
Prime-index labels with $\ell\ge5$ distinguish all three values and have
entropy $\log3$, which is larger. Thus convex combinations of these
coordinatewise bounds and only the two aggregate size budgets cannot push the
constant below $C_0$. This is a barrier for that relaxation, not a lower bound
for an actual arithmetic fiber.
